\section{Finite certificates for the whole numerical experiment}\label{sec:cert46}
The Hankel normal form uses all suffixes. We now give an equivalent
finite certificate whose size is polynomial in a \emph{displayed}
horizon and proposed widths. We then treat the actual optimal error,
not the sum of local compatibility defects. Throughout this section
$X$ and $\calA$ are nonempty, $|\calA|=m$, $D\ge1$, and
$K_0,\ldots,K_N\ge1$. The held binary answer is a separate two-label
cut. A horizon is supplied in unary for statements about input size.
Rational numbers have binary encodings; a real algebraic coefficient
is supplied by a polynomial and a rational isolating interval.
These conventions do not give an oracle for arbitrary real numbers.

In this section $T_{t,a}$ labels the edge from cut $t$ to cut $t+1$,
a reindexing of Definition~\ref{def:machine}. Write $E_x$ for an initialization row, $T_{t,a}$ for the
$K_t$-by-$K_{t+1}$ row-stochastic transition, and $d_j\in[-1,1]^{K_N}$
for the terminal mean column. Define the residual mean
\begin{equation}\label{eq:res46}
 R_{x,w,j}=E_xT_{0,a_1}\cdots T_{N-1,a_N}d_j
             -\rho\ip{e_j}{U_{a_N}\cdots U_{a_1}x}.
\end{equation}
Thus the actual worst-word binary-TV error is
$e(\mathcal M)=\frac12\max_{x,w,j}|R_{x,w,j}|$.
In particular, centering the numerical response does not turn a
bounded stochastic decoder into an unrestricted signed readout.

Set $n_t=K_t+D$ and form rectangular block matrices and vectors
\begin{equation}\label{eq:aug46}
 A_{t,a}=\operatorname{diag}(T_{t,a},U_a^{\mathsf T}),\quad
 u_x=(E_x,-\rho x^{\mathsf T}),\quad
 v_j=\binom{d_j}{e_j}.
\end{equation}
The order and transpose in this definition are essential for
noncommuting commands. Equation~\eqref{eq:res46} equals $u_xA_wv_j$.

\begin{theorem}[A finite all-word certificate]\label{thm:gram46}
Define
\begin{equation}\label{eq:gram46}
 G_0=\sum_{x\in X}u_x^{\mathsf T}u_x,\qquad
 G_{t+1}=\sum_{a\in\calA}A_{t,a}^{\mathsf T}G_tA_{t,a},\qquad
 S_2=\sum_{j=1}^D v_j^{\mathsf T}G_Nv_j.
\end{equation}
Then
\[
 S_2=\sum_{x\in X}\sum_{w\in\calA^N}\sum_{j=1}^D R_{x,w,j}^2.
\]
Consequently a legal proposed machine is exact on every word if and
only if $S_2=0$. For rational data this can be checked by exact
arithmetic in polynomial time in the displayed data. If $S_2>0$, a
seed, word and query with nonzero residual can also be found in
polynomial time, without listing all words.

Exact feasibility of a prescribed profile for an algebraic
orthogonal-word target has an existential real formula of polynomial
size in $N,m,D,|X|$ and the explicitly displayed widths. Its equations
may all be taken to have degree at most three after auxiliary arithmetic
variables are introduced. Feasibility has a real algebraic witness
whenever it has a real witness.
\end{theorem}
\begin{proof}
Induction in \eqref{eq:gram46} gives
$G_t=\sum_{x,|u|=t}(u_xA_u)^{\mathsf T}(u_xA_u)$.
Pairing with $v_j$ proves the identity. Every summand is a real square,
so zero means equality of all numerical responses, not merely equality
on average. No assumption about the ranks of reachable distributions
or the observability of individual states is used.

For a witness, compute backwards
\[
 C_N=\sum_jv_jv_j^{\mathsf T},\qquad
 C_t=\sum_a A_{t,a}C_{t+1}A_{t,a}^{\mathsf T}.
\]
Select $x$ with $u_xC_0u_x^{\mathsf T}>0$. Given a current row $u$ at
cut $t$ with $uC_tu^{\mathsf T}>0$, at least one command satisfies
$uA_{t,a}C_{t+1}A_{t,a}^{\mathsf T}u^{\mathsf T}>0$.
Append it and replace $u$ by $uA_{t,a}$. At the last cut choose $j$
with $uv_j\ne0$. The matrices are positive semidefinite by their
recursions, so this procedure cannot encounter a dead end.
There are $O(Nm)$ matrix operations at dimensions at most
$\max_t n_t$. For rational input, common-denominator bounds on these
finite matrix products also give polynomial intermediate bit length.
Exact algebraic sign testing gives the corresponding effective
statement for algebraic inputs; no unspecified real-number comparison
is invoked.

For existence, treat $E,T,d,G$ as unknowns. Impose nonnegativity and
unit row sums on $E,T$, the bounds on $d$, the recursions
\eqref{eq:gram46}, and $S_2=0$. Each recursion entry has degree at most
three, and its number of terms is polynomial in the matrix sizes.
The Gram variables need no independent positivity constraint: the
initial condition and recurrence already force the displayed sum of
squares. If target coefficients are algebraic rather than rational,
include their defining polynomials and isolating inequalities. Arithmetic
circuit variables reduce larger degrees to quadratic equations.
Elimination over real closed fields decides existence and permits
algebraic sample points. This last conclusion is the usual
real-algebraic sampling principle \cite{BPR}, not a polynomial-time optimizer.
\end{proof}

The linear-algebraic equivalence mechanism is classical: compare
Tzeng's equivalence algorithm \cite{Tzeng92} and the weighted-automata
Hankel and Gram constructions of Balle--Panangaden--Precup
\cite{BPP15}. The assertion here makes the finite clock, stochastic
constraints, all seed initializations, and bounded query columns
simultaneous. It must not be read as a new discovery of polynomial-time
weighted-automata equivalence. Also, $S_2\approx0$ in floating point is
\emph{not} the exact certificate $S_2=0$.

\subsection{A hierarchy for the actual optimal error}
Let $Q=|X|Dm^N$. For each even integer $p\ge2$, define
\begin{equation}\label{eq:momnorm46}
 L_p(\mathcal M)=\frac12
       \left(Q^{-1}\sum_{x,w,j}R_{x,w,j}^{p}\right)^{1/p},\qquad
 \ell_p(K_\bullet)=\min_{\mathcal M}L_p(\mathcal M),
\end{equation}
where every minimization uses the \emph{same} product of stochastic row
simplices and bounded decoder cubes for the specified profile. Put
$\mathcal E(K_\bullet)=\min_{\mathcal M}e(\mathcal M)$.

\begin{theorem}[Moment certificates and the error frontier]\label{thm:moments46}
All the preceding minima are attained. For every even $p\ge2$,
\begin{equation}\label{eq:frontier46}
 \ell_p(K_\bullet)\le\mathcal E(K_\bullet)
                 \le Q^{1/p}\ell_p(K_\bullet).
\end{equation}
The numbers $\ell_p$ increase along the even integers and converge to
$\mathcal E$. The running minima of the right sides of
\eqref{eq:frontier46} decrease to the same value. In particular,
$\ell_2=0$ if and only if the profile admits an exact realization.
For each fixed $p$, its moment sum is computable by a finite homogeneous
polynomial recurrence without enumerating words. For fixed
$\max_t(K_t+D)$, this recurrence has size polynomial in $p,N,m,D,|X|$.

For algebraic target data the actual optimum $\mathcal E$ and a
minimizing machine are algebraic and effectively obtainable. For a
finite prescribed width cap, the resulting threshold values give the
exact finite-horizon error/width frontier. This is an effective finite
optimization statement, not a closed formula or a polynomial-time
algorithm for growing width.
\end{theorem}
\begin{proof}
The parameter set is nonempty and compact, and the objectives are
continuous. The inequalities for a vector of $Q$ numbers are
\[
 \left(Q^{-1}\sum |r_i|^p\right)^{1/p}
 \le\max_i|r_i|\le
 Q^{1/p}\left(Q^{-1}\sum |r_i|^p\right)^{1/p}.
\]
The first inequality survives minimization. For the second, use a
minimizer of $L_p$ as a candidate for $\mathcal E$; no exchange of a
maximum over words with a minimum over machines is made. Normalized
finite-measure $L^p$ norms increase with $p$, as do their minima.
Since $Q^{1/p}\to1$, the two sequences converge as asserted, including
the case of a zero optimum.

The recurrence is particularly transparent in polynomial form. Let
$z$ be a row variable of the appropriate dimension and set
\begin{equation}\label{eq:sympow46}
 F_N(z)=\sum_j(zv_j)^p,\qquad
 F_t(z)=\sum_a F_{t+1}(zA_{t,a}).
\end{equation}
Then $\sum_xF_0(u_x)$ is the moment sum. A homogeneous degree-$p$
polynomial in $n_t$ variables has $\binom{n_t+p-1}{p}$ coefficients.
Repeated polynomial multiplication and linear substitution compute
these coefficients in polynomially many arithmetic operations when
$n_t$ is fixed. This is the symmetric-power form of the familiar tensor
product of weighted representations. It does not imply a bound
polynomial in both $p$ and unbounded $n_t$.

For the last assertion one may instead introduce all finitely many
word inequalities $-2e\le R_{x,w,j}\le2e$ and minimize $e\in[0,1]$.
Projection of this compact semialgebraic feasible set gives its
semialgebraic attainable-error set. Its left endpoint is algebraic,
and real-algebraic sampling gives a machine above that endpoint.
The horizon and word set are finite; the number of those inequalities
can be exponential. Standard real quantifier elimination proves the
claim, with no missing infinite-horizon limit. Apply it successively
to the finitely many allowed integer profiles or to padded constant
caps. The held two-label output cut is then included in the reported
peak.\qedhere
\end{proof}

For a \emph{given} machine and tolerance factor $1+\delta$, choosing
an even $p\ge\log Q/\log(1+\delta)$ gives lower and upper error
certificates with ratio at most $1+\delta$. This statement concerns
checking the supplied machine, not finding a global moment minimizer.
A certificate for a globally optimized bound must certify that
optimization as well. A numerical local minimum has no such status.
For fixed dimensions, exact rational coefficient growth in
\eqref{eq:sympow46} is bounded by polynomially many input bit lengths,
products and multinomial factors; certified root intervals can replace
irrational $p$th roots in the output. The implementation below uses
exact degree-two and degree-four recurrences, not an unreported
large-degree optimization.

\subsection{Finite precision without changing boundary widths}
For an integer $L\ge1$, let $\mathcal G_L$ contain all machines whose
initial and transition probabilities, and positive-answer probabilities,
are integral multiples of $1/L$. Let
$e_L=\min_{\mathcal M\in\mathcal G_L}e(\mathcal M)$ and set
\begin{equation}\label{eq:gridcost46}
 C(K_\bullet)=(K_0-1)+\sum_{t=1}^N(K_t-1)+1.
\end{equation}
These finite sets are nonempty even when some cuts have only one label.

\begin{theorem}[A finite two-sided rational certificate]\label{thm:grid46}
For every finite binary target, without any orthogonality assumption,
\begin{equation}\label{eq:grid46}
 \max\{0,e_L-C(K_\bullet)/L\}
       \le\mathcal E(K_\bullet)\le e_L.
\end{equation}
Consequently dyadic machines with the same available boundary profile
attain every strictly larger error than $\mathcal E$. The theorem
does not assert a dyadic or rational exact optimizer at the boundary.
If two target probability tables differ by at most $\eta$ uniformly,
their optimal errors at the same profile differ by at most $\eta$.
Thus a strict positive feasibility or infeasibility margin persists
under sufficiently small target perturbations.
\end{theorem}
\begin{proof}
Round down the first $k-1$ coordinates of any probability row of size
$k$ to multiples of $1/L$, and give the remaining mass to the last
coordinate. The row stays stochastic and its TV change is at most
$(k-1)/L$. Round a terminal positive-answer probability to within $1/L$.
Couple the two machines at initialization and then at each common
command, or telescope the products by TV contraction. For every word
the total output change is bounded by \eqref{eq:gridcost46} divided by
$L$. Applying this construction to a minimizing real machine gives
$e_L\le\mathcal E+C/L$; inclusion of the grid in the full parameter
set gives $\mathcal E\le e_L$. Both use the same labels; no scratch
configuration is hidden in that assertion. Choosing $L=2^b$ sufficiently
large proves the strict-margin assertion. Finally, for any fixed
machine its errors against two targets differ by at most $\eta$ by the
triangle inequality. Minimize in both directions.
\end{proof}

An exhaustive grid search is a terminating certificate algorithm at
fixed $L$, not an efficient proposal for large instances. If the target
is supplied by certified intervals of radius $\eta$, the two endpoints
in \eqref{eq:grid46} can each be enlarged by $\eta$. Dyadic row sampling
uses $b$ fresh fair bits at denominator $2^b$; index, table and scratch
storage still have to be charged in an implementation. The separate
uniform compiler below supplies its own workspace accounting. The
rational approximation statement must not be strengthened to exact
rational factorization: such field restrictions genuinely matter in
nonnegative realization \cite{Shitov18}.

\begin{proposition}[An exact calibrated finite frontier]\label{prop:smallfront46}
Take $N=0$, the four planar seeds $\{\pm e_1,\pm e_2\}$, two coordinate
queries and $0<\rho\le1/2$ (this wider signal range is used only
for the present finite example). If $k$ is the available initialization
alphabet, the optimal binary-TV errors are
\begin{equation}\label{eq:smallfront46}
 \mathcal E(1)=\rho/2,\qquad
 \mathcal E(2)=\rho/4,\qquad
 \mathcal E(k)=0\quad(k\ge3).
\end{equation}
The same formula holds for an arbitrary horizon consisting only of
identity commands, with $k=\min_t K_t$. In all cases the final held
answer still costs two labels.
\end{proposition}
\begin{proof}
A one-state decoder is one point of the mean square. Approximating
$\pm\rho e_1,\pm\rho e_2$ by a common point has minimum uniform
coordinate error $\rho$, attained at zero. For two states every seed
mean lies on one affine line $n\cdot y=c$, with $n\ne0$; a degenerate
segment may be embedded in any such line. If all coordinate errors
are at most $d$, duality of $\ell_\infty$ and $\ell_1$ gives
\[
 d\|n\|_1\ge
 \max_{y\in\{\pm\rho e_1,\pm\rho e_2\}}|n\cdot y-c|
 =\rho\|n\|_\infty+|c|\ge\tfrac\rho2\|n\|_1.
\]
The diagonal segment with endpoints $\pm(\rho/2,\rho/2)$ attains
coordinate error $\rho/2$: encode $e_1,e_2$ at its positive endpoint
and their negatives at its negative endpoint. Divide mean errors by
two to obtain binary TV. Three legal mean vectors
$(2\rho,0),(-\rho,2\rho),(-\rho,-2\rho)$ enclose all four targets,
so their barycentric initialization gives exact responses.
For identity-only words, condition at a narrowest cut to obtain the
same one-cut lower bound. Initialize the appropriate construction and
carry its label through all cuts, padding unused labels as necessary,
to obtain the upper bound.
\end{proof}

\subsection{A complexity boundary with a different output convention}
The following comparison concerns one \emph{categorical} terminal
output, not a selected member of a family of binary queries.
Its width is the pre-output memory alphabet; counting the held output
would separately add its alphabet size.

\begin{proposition}[Normalized factorization and the real field]\label{prop:hard46}
Given an explicitly listed rational row-stochastic matrix $P$, deciding
whether it has a stochastic factorization $P=ED$ of inner dimension at
most $k$ is $\exists\R$-complete. It is already a zero-command numerical
realization problem. Rational target data need not have a rational
exact optimizer at the minimum real inner dimension.
\end{proposition}
\begin{proof}
Membership follows from the nonnegative variables, row-sum equations
and quadratic equations $P=ED$. For hardness, remove the zero rows of
a nonzero nonnegative rational matrix $A$ and divide each remaining row
by its positive row sum, giving $P$. Positive diagonal row scaling and
removal of zero rows preserve nonnegative rank. If $P=WH$ is a
nonnegative factorization, put $h_i=\sum_jH_{ij}$. For $h_i>0$ set
$D_{ij}=H_{ij}/h_i$ and $E_{xi}=W_{xi}h_i$. If $h_i=0$, replace
that unused decoder row by any probability row and set its $E$ column
to zero. Then $D\mathbf1=\mathbf1$ and
$E\mathbf1=P\mathbf1=\mathbf1$. Thus stochastic inner dimension
is exactly nonnegative rank, over either the real or rational field.
Shitov's Theorem~2 gives $\exists\R$-hardness and his Corollary~3
gives rational matrices with different minimum ranks over the rational
and real algebraic fields \cite{Shitov18}. Row normalization transfers
both statements as just shown. Zero input matrices can be treated
separately and do not affect the reduction.
\end{proof}

This transfers a known complexity theorem; it is not a new hardness
proof for nonnegative rank. Nor does it establish the same hardness for
the fixed-dimensional orthogonal binary-query subclass. In that
subclass Theorem~\ref{thm:gram46} supplies the compact exact feasibility
formula, while the inherited arithmetic theorems determine specific
asymptotic widths. These are different assertions.
